\scanpage{55}
\mspage{23}
avec la structure de stratification cylindrique \inserted{\uline{ordonnée}}
locale. \deleted{Par exemple} Une autre façon de
décrire les $\Gamma(d)=\Sigma^{<}(d+1)$, c'est de
prendre la puiss. \deleted{\unsure{sym.}} \inserted{cart. relative} $d$.ième de $\Gamma_0$\deleted{\unsure{$(\varnothing)$}}\inserted{$=\Sigma^{<}(1)$}
sur $\Sigma\ (=\Gamma(\varnothing))$, dont on enlève les images
des puiss. cart. relatives \deleted{inférieures} \inserted{$\Gamma'(d')$ des $d'$ième ($d'<d$)} par les applic\unsure{ns}
$\unsure{\tfrac12}$ simpliciales entre celles-ci. Cela donne à priori
les \deleted{\unsure{ens tot.}} suites ordonnées de $d$ germes de
strates de codim $1$ \deleted{(\unsure{mutuellement}} \inserted{en un pt,} \unsure{distinctes|distincts}), \deleted{soit}
$\Gamma'(d)$. On a alors \inserted{$\unsure{\Gamma^{\deleted{'}}(d)|\Gamma(d)}\subset\Gamma'(d)$},
$\Gamma(d)$ est une partie à la fois ouverte et
\deleted{fermée} de $\Gamma'(d)$, et \deleted{$\Gamma'(d)\simeq$}
\[
(35)\qquad \Gamma'(d)\simeq\unsure{\mathfrak S_d|\mathcal S_d}\times\Gamma(d).
\]
Pour avoir \deleted{la}\inserted{une} structure d'ordre \inserted{total} relative de
\deleted{\unsure{$\Gamma'(1)$ sur}} $\Gamma_0=\Sigma^{<}(1)$ sur $\Sigma_{\deleted{\unsure{*}}}=\Sigma^{<}(0)$, il \deleted{suffit}\inserted{\unsure{uniquement}}
de se donner, \deleted{dans $\Gamma'_1(0)=\unsure{\Sigma^{<}(0)}$, $\Gamma'_{\unsure{1+1}}$} dans
$\Gamma_0\times_\Sigma\Gamma_0\setminus\text{diagonale}$, une partie \inserted{(i.e.} ouverte et fermée)
$\Gamma_1$, telle que $\Gamma_1\cup{}^s\Gamma_1=\Gamma_0\times\Gamma_0\setminus\text{diag}$, $\Gamma_1\cap{}^s\Gamma_1=\varnothing$
i.e. une section de $\Gamma_0\times_\Sigma\Gamma_0\setminus\text{diag}$ au dessus
de $\operatorname{Sym}^2_\Sigma(\Gamma_0)^*$ (``choix d'un premier élément parmi deux
éléments distincts), avec la condition de transitivité
\margin{\deleted{il s'agit de}}
\margin{$\Gamma_1$ est le graphe de la relation $x<y$.}
\margin{ou une section \deleted{de $\Gamma_0\times_\Sigma\Gamma_0\setminus\text{diagonale}$}\deleted{de $\operatorname{Sym}^2(\unsure{\Sigma\times_\Sigma\Gamma_0})$}}
\deleted{$\unsure{\text{que }\Gamma_1}$}
\diagram{\begin{tikzcd}[column sep=large,row sep=small]
(\Gamma_1\times_\Sigma\Gamma_0)\cap(\Gamma_0\times_\Sigma\Gamma_1)
 \arrow[r,"{\operatorname{pr}_{13}}"] \arrow[d,phantom,"\cap" description]
 & \Gamma_0\times\Gamma_0 \arrow[d,phantom,"\cup" description] \\
\Gamma_0\times\Gamma_0\times\Gamma_0 & \Gamma_1
\end{tikzcd}}
(36) \quad se factorise par $\Gamma_1$
\mspage{54}
